\documentclass[a4paper,12pt]{article}
\usepackage[T2A]{fontenc}
\usepackage[utf8]{inputenc}
\usepackage[russian]{babel}
\usepackage[a4paper,margin=23mm]{geometry}
\usepackage{amsmath,amssymb,booktabs,tabularx,xcolor,tikz,fancyhdr}
\usepackage[colorlinks=true,linkcolor=teal!65!black]{hyperref}
\definecolor{accent}{HTML}{006666}
\setlength{\parindent}{0pt}
\setlength{\parskip}{5pt}
\setlength{\headheight}{15pt}
\pagestyle{fancy}
\fancyhf{}
\fancyhead[L]{\small Сравнения по модулю}
\fancyhead[R]{\small Ксения Клыкова}
\fancyfoot[C]{\thepage}
\renewcommand{\headrulewidth}{0.3pt}
\newcommand{\block}[1]{\par\medskip\textbf{#1}\par\nopagebreak[3]}
\newcommand{\err}[1]{\block{\textcolor{red!75!black}{Ошибка:} #1}}
\newcommand{\idea}{\block{\textcolor{blue!60!black}{Ключевая идея}}}
\newcommand{\result}[1]{\par\textcolor{accent}{\textbf{Итог:}} #1}
\newcommand{\task}[1]{\clearpage\phantomsection\label{task:#1}\section*{\textcolor{accent}{Задание № #1}}}
\newcommand{\practice}{\block{Задача на закрепление}}
\begin{document}
\begin{titlepage}
\centering
\vspace*{25mm}
{\Large\color{accent} МАТЕМАТИКА\par}
\vspace{15mm}
{\huge Ревью домашней работы\par}
\vspace{9mm}
{\LARGE\color{accent} Сравнения по модулю\par}
\vfill
{\large Ксения Клыкова\par}
\vspace{9mm}
28 ноября 2026 г.\par
\vfill
{\small Лёвин Артём Александрович\\эксклюзивно для Ксении Клыковой\par}
\vspace{14mm}
\begin{tikzpicture}
\draw[accent,line width=.7pt] (0,0) .. controls (2,.7) and (4,-.7) .. (6,0) .. controls (8,.7) and (10,-.7) .. (12,0);
\end{tikzpicture}
\end{titlepage}
\setcounter{page}{1}
\section*{\textcolor{accent}{Сводная таблица}}
\small
\renewcommand{\arraystretch}{1.45}
\begin{tabularx}{\linewidth}{@{}c >{\raggedright\arraybackslash}p{.20\linewidth} >{\raggedright\arraybackslash}X >{\raggedright\arraybackslash}p{.20\linewidth} >{\raggedright\arraybackslash}p{.21\linewidth}@{}}
\toprule
№ & Тема & Суть ошибки & Ваш ответ & Верный ответ\\
\midrule
\hyperref[task:2]{2} & Остатки & Ошибка при замене 109 & В вычислении $-1$ & $1$\\
\hyperref[task:5]{5} & Многочлены & Равенство вместо сравнения & $5$ & $5$\\
\hyperref[task:6]{6} & Линейное сравнение & Неверный модуль при перечислении & $3,7,11,15,19,23$ (mod 4) & $3,7,11,15,19,23$ (mod 24)\\
\hyperref[task:7]{7} & Степени & Не учтена чётность показателя & $1; -1$ & $1;1$\\
\hyperref[task:8]{8} & Квадраты & Пропущен остаток 0 & Нет & Нет; полный перебор\\
\hyperref[task:10]{10} & Делимость & Неверный объект делимости в выводе & Общая часть верна; вывод про $n$ & $11\mid(4^{10001}+7^{10001})$\\
\bottomrule
\end{tabularx}
\normalsize
\clearpage
\thispagestyle{empty}
\vspace*{\fill}
\begin{center}{\Huge\color{accent} Разбор ошибок}\end{center}
\vspace*{\fill}
\task{2}
\block{Текст задания}
Найдите остаток от деления $109\cdot46+124$ на 11.
\block{Ваше решение}
\[
109\equiv-2\pmod{11},\quad 46\equiv2\pmod{11},\quad124\equiv3\pmod{11}.
\]
В обосновании записано $109=11\cdot10-2$. Далее: $-2\cdot2+3=-1$.
\block{Ваш ответ}
В последней строке вычислений получено $-1$. Это значение не является стандартным остатком по модулю 11.
\err{Неверное разложение числа 109.}
\block{Где именно возникает ошибка}
Вы правильно выбрали способ: заменить каждое число удобным представителем его класса остатков. Первый неверный шаг --- $109=110-2$: правая часть равна 108. Последующие сравнения для 46 и 124 верны.
\idea
По определению $a\equiv b\pmod m$, если разность $a-b$ кратна $m$. Для стандартного остатка дополнительно требуется $0\le r<m$.
\block{Исправление и полное решение}
Поскольку $109=110-1$, получаем $109\equiv-1\pmod{11}$. Проверяем остальные замены: $46=44+2$, $124=121+3$.
Сравнения можно умножать и складывать при одинаковом модуле:
\[
109\cdot46+124\equiv(-1)\cdot2+3=1\pmod{11}.
\]
Число 1 находится в диапазоне от 0 до 10.
\block{Проверка}
$109\cdot46+124=5138=11\cdot467+1$.
\result{Стандартный остаток равен 1.}
\practice
Найдите стандартный остаток от деления $131\cdot38+95$ на 11.
\task{5}
\block{Текст задания}
Известно, что $n\equiv-2\pmod9$. Найдите остаток от деления $n^3-2n^2+5n+4$ на 9.
\block{Ваше решение и ответ}
Вы заменили $n$ числом $-2$, записав равенство многочлена с числовым выражением, получили $-22$, затем $-22\equiv5\pmod9$. Ответ: 5.
\err{Подмена сравнения обычным равенством.}
\block{Где именно возникает ошибка}
Условие $n\equiv-2\pmod9$ прочитано правильно. Следующий переход, записанный со знаком равенства, не следует из условия: $n$ может быть любым числом вида $9k-2$, где $k$ целое. Например, $n=7$ также удовлетворяет условию.
\idea
При замене целых чисел сравнимыми по модулю числами сумма, произведение и натуральные степени сохраняют сравнимость. Значения многочлена при разных аргументах могут различаться на кратное модуля.
\block{Исправление и полное решение}
\begin{align*}
n^3-2n^2+5n+4&\equiv(-2)^3-2(-2)^2+5(-2)+4\pmod9\\
&\equiv-8-8-10+4=-22\pmod9.
\end{align*}
Теперь $-22=9\cdot(-3)+5$, следовательно, $-22\equiv5\pmod9$. Число 5 является стандартным остатком.
\block{Проверка}
Возьмём $n=7$. Тогда $343-98+35+4=284=9\cdot31+5$. При $n=-2$ значение равно $-22$. Значения различны, но оба имеют остаток 5.
\result{Численный ответ верен; при замене аргумента нужен знак $\equiv$.}
\practice
Известно, что $n\equiv-3\pmod8$. Найдите стандартный остаток от деления $n^3-2n^2+4n+5$ на 8. Запишите преобразования с корректными знаками.
\task{6}
\block{Текст задания}
Известно, что $6x\equiv18\pmod{24}$. Запишите все возможные стандартные остатки $x$ по модулю 24.
\block{Ваше решение и ответ}
Вы получили $x\equiv3\pmod4$ и перечислили числа $3,7,11,15,19,23$, подписав их модулем 4.
\err{Модуль в итоговом перечислении.}
\block{Где именно возникает ошибка}
Сокращение исходного сравнения выполнено верно. Все шесть чисел тоже найдены верно. Ошибка возникает при указании модуля 4 после списка: стандартные остатки по модулю 4 могут быть только 0, 1, 2, 3.
\idea
Сокращение сравнения меняет модуль. Если $da\equiv db\pmod{dm}$, где $d,m$ положительные целые, то $a\equiv b\pmod m$. Однако исходный вопрос может требовать перечислить представители по первоначальному модулю.
\block{Исправление и полное решение}
По определению $24\mid(6x-18)$, то есть $6(x-3)=24k$ при некотором целом $k$. Делим на 6: $x-3=4k$, откуда $x=3+4k$. Для стандартных остатков по модулю 24 нужно $0\le3+4k\le23$, поэтому $k=0,1,2,3,4,5$.
\[
\begin{array}{c|rrrrrr}
k&0&1&2&3&4&5\\\hline
x&3&7&11&15&19&23
\end{array}
\]
\block{Проверка}
Для каждого такого $x$ имеем $6x=18+24k$, поэтому сравнение выполняется. Следующее число 27 уже выходит за диапазон стандартных остатков по модулю 24.
\result{$3,7,11,15,19,23$ --- все стандартные остатки по модулю 24.}
\practice
Решите сравнение $8x\equiv24\pmod{40}$ и перечислите все стандартные остатки $x$ по модулю 40.
\task{7}
\block{Текст задания}
Найдите остатки числа $10^{2026}$ при делении на 9 и на 11.
\block{Ваше решение и ответ}
Записано $10^{2026}\equiv1\pmod9$ и $10^{2026}\equiv-1\pmod{11}$.
\err{Потеря чётности показателя степени.}
\block{Где именно возникает ошибка}
Первое сравнение верно. Во втором сравнении итог $-1$ неверен: основание 10 действительно сравнимо с $-1$ по модулю 11, но при возведении в чётную степень знак становится положительным. Промежуточные шаги в работе не записаны, поэтому точный ход вычислений установить нельзя.
\idea
Если $a\equiv b\pmod m$, то $a^s\equiv b^s\pmod m$ для положительного целого $s$. При чётном $s$ число $(-1)^s$ равно 1; при нечётном равно $-1$.
\block{Исправление и полное решение}
По модулю 9: $10=9+1$, значит $10^{2026}\equiv1^{2026}=1\pmod9$.

По модулю 11: $10=11-1$, значит
\[
10^{2026}\equiv(-1)^{2026}=\bigl((-1)^2\bigr)^{1013}=1\pmod{11}.
\]
В обоих случаях 1 входит в диапазон стандартных остатков.
\block{Проверка}
$10^2=100=99+1$. Поэтому $10^{2026}=100^{1013}$ сравнимо с 1 по модулю 9 и по модулю 11, так как 99 кратно каждому из них.
\result{Оба остатка равны 1.}
\practice
Найдите стандартные остатки числа $10^{2027}$ при делении на 9 и на 11. Обоснуйте каждый ответ.
\task{8}
\block{Текст задания}
Может ли квадрат целого числа оканчиваться на 8? Обоснуйте ответ перебором возможных остатков.
\block{Ваше решение и ответ}
Вы проверили квадраты чисел от 1 до 9. Последние цифры: $1,4,9,6,5,6,9,4,1$. Ответ: нет.
\err{Неполный перебор классов остатков.}
\block{Где именно возникает ошибка}
Все перечисленные квадраты и их последние цифры верны. Однако в переборе отсутствует класс чисел, кратных 10, то есть остаток 0. Условие требует рассмотреть все возможные остатки, включая нулевой.
\idea
Последняя цифра целого неотрицательного числа --- его стандартный остаток по модулю 10. Для квадрата любого целого $n$ можно использовать остаток самого $n$: если $n\equiv r\pmod{10}$, то $n^2\equiv r^2\pmod{10}$.
\block{Исправление и полное решение}
Добавим случай $r=0$, когда $r^2=0$. Полный перебор:
\[
\begin{array}{c|rrrrrrrrrr}
r&0&1&2&3&4&5&6&7&8&9\\\hline
r^2\bmod10&0&1&4&9&6&5&6&9&4&1
\end{array}
\]
Множество возможных последних цифр квадрата равно $\{0,1,4,5,6,9\}$. Цифры 8 в нём нет.
\block{Проверка полноты}
Каждое целое число, включая отрицательные, имеет ровно один стандартный остаток от 0 до 9 по модулю 10. Все десять случаев учтены, поэтому перебор исчерпывающий.
\result{Квадрат целого числа не может оканчиваться на 8. Ответ верен; доказательство дополнено.}
\practice
Может ли квадрат целого числа оканчиваться на 7? Обоснуйте ответ полным перебором остатков по модулю 10.
\task{10}
\block{Текст задания}
Докажите, что $4^n+7^n$ делится на 11 при каждом положительном нечётном целом $n$. Примените результат к $n=10^4+1$.
\block{Ваше решение и ответ}
Вы правильно использовали $4\equiv-7\pmod{11}$ и нечётность $n$, получив $-7^n+7^n\equiv0\pmod{11}$. В заключении определили нечётность $n=10^4+1$, но записали вывод с обозначением $11\mid n\equiv0$.
\err{В заключении перепутаны показатель степени и проверяемое выражение.}
\block{Где именно возникает ошибка}
Общее доказательство верно. Последний корректный шаг --- установление нечётности $10^4+1$. После него нужно применить доказанное утверждение к сумме степеней. Делимость самого показателя на 11 из нечётности не следует.
\idea
Запись $m\mid A$ означает, что $A$ кратно $m$, и равносильна $A\equiv0\pmod m$. В этой задаче $A=4^n+7^n$.
\block{Исправление и полное решение}
Так как $4-(-7)=11$, имеем $4\equiv-7\pmod{11}$. Возводим обе части в натуральную степень $n$. Для нечётного $n$ получаем $(-7)^n=-7^n$. Тогда
\[
4^n+7^n\equiv(-7)^n+7^n=-7^n+7^n=0\pmod{11}.
\]
Теперь $10^4+1=10001$ --- положительное нечётное число, поэтому доказанное утверждение применимо:
\[
4^{10001}+7^{10001}\equiv0\pmod{11}.
\]
\block{Проверка}
Сам показатель $10001=11\cdot909+2$ имеет остаток 2. Сумма степеней имеет остаток 0 по доказанному правилу. Например, при $n=3$ сумма $64+343=407=11\cdot37$.
\result{$11\mid(4^{10001}+7^{10001})$.}
\practice
Докажите, что $5^n+8^n$ делится на 13 при каждом положительном нечётном целом $n$. Примените результат к $n=10^3+1$.
\vfill
\begin{center}\small Лёвин Артём Александрович эксклюзивно для Ксении Клыковой\end{center}
\end{document}
